\subsubsection{Fundamental Solutions}\label{subsec:fundamental-solutions}

\begin{Definition}
\label{def:fundamental-solution}
Let $L(D)$ be a differential operator with constant coefficients. Then a \textbf{fundamental solution} for $L$ is a distribution $G\in\mathcal{D}^{\prime}(\mathbb{R}^m)$ satisfying the equation $L(D)G=\delta$.
\end{Definition}

Of course, fundamental solutions are unique only up to a solution of the homogeneous equation $L(D)u=0$; in choosing a specific fundamental solution one often selects the one with the "nicest" behavior at infinity. The significance of the fundamental solution lies in the fact that
\begin{equation}\label{eq:fundamental-solution-convolution-identity}
L(D)(G*f)=(L(D)G)*f=\delta*f=f,
\end{equation}
provided that the convolution $G*f$ is defined. If, for example, $f$ has compact support, then $G*f$ is a solution of the equation $L(D)u=f$.

The construction of fundamental solutions for general operators with constant coefficients is quite complicated, and we shall limit our discussion to some important examples.

\begin{Example}{Ordinary differential equations}
\label{ex:ode-fundamental-solution}
We seek a solution to the ODE
\begin{equation}\label{eq:ode-fundamental-solution-equation}
a_nG^{(n)}(x)+\cdots+a_0G(x)=\delta(x).
\end{equation}
For both positive and negative $x$, $G$ must agree with a solution of the homogeneous ODE. That is, if $u_1(x)$, $\ldots$, $u_n(x)$ are a complete set of solutions for the homogeneous ODE, then we must have
\begin{equation}\label{eq:ode-fundamental-solution-piecewise-basis}
G(x)=\begin{cases}\alpha_1u_1(x)+\cdots+\alpha_nu_n(x)&x>0\\\beta_1u_1(x)+\cdots+\beta_nu_n(x)&x<0.\end{cases}
\end{equation}
We can now satisfy \eqref{eq:ode-fundamental-solution-equation} by requiring that all derivatives of $G$ up to the $(n-2)$nd are continuous at 0, but the $(n-1)$st derivative has a jump of magnitude $1/a_n$. With $\gamma_i=\alpha_i-\beta_i$, this yields the system
\begin{equation}\label{eq:ode-fundamental-solution-jump-system}
\begin{aligned}
\gamma_1u_1(0)+\cdots+\gamma_nu_n(0)&=0,\\
\gamma_1u_1'(0)+\cdots+\gamma_nu_n'(0)&=0,\\
&\vdots\\
\gamma_1u_1^{(n-2)}(0)+\cdots+\gamma_nu_n^{(n-2)}(0)&=0,\\
\gamma_1u_1^{(n-1)}(0)+\cdots+\gamma_nu_n^{(n-1)}(0)&=\dfrac1{a_n}.
\end{aligned}
\end{equation}
The determinant of this system is the Wronskian of the complete set of solutions $u_i$ and is hence nonzero.
\end{Example}

\Needspace{4\baselineskip}
\begin{Example}{Laplace\textquotesingle s equation}
\label{ex:laplace-fundamental-solution}
We now seek a solution of the equation
\begin{equation}\label{eq:laplace-fundamental-solution-equation}
\Delta G(\bm{x})=\delta(\bm{x})
\end{equation}
on $\mathbb{R}^m$. Of course this makes $G$ a solution of the homogeneous Laplace equation except at the origin. Moreover, since $\delta$ is radially symmetric, it is natural to seek a radially symmetric $G$. For radially symmetric functions, Laplace\textquotesingle s equation reduces to
\begin{equation}\label{eq:laplace-radial-fundamental-ode}
G''(r)+\dfrac{m-1}rG'(r)=0,\quad r>0,
\end{equation}
and we obtain the solution
\begin{equation}\label{eq:laplace-radial-homogeneous-solutions}
G(r)=\begin{cases}c_1+c_2r^{2-m}&m\geq3\\c_1+c_2\ln r&m=2.\end{cases}
\end{equation}
We can now satisfy \eqref{eq:laplace-fundamental-solution-equation} by an appropriate choice of $c_2$. For $m=3$, we did this calculation in Example~\ref{ex:inverse-radius-laplacian} of the previous section; the result was $c_2=-1/4\pi$. The general result is obtained in an analogous fashion; one finds the fundamental solutions
\begin{equation}\label{eq:laplace-fundamental-solution-all-dimensions}
G(\bm{x})=\begin{cases}-r^{2-m}/(m-2)\Omega_m,&m\geq3\\\ln r/2\pi&m=2.\end{cases}
\end{equation}
Here $\Omega_m=2\pi^{m/2}/\Gamma(m/2)$ is the surface area of the unit sphere (cf. Problem~\ref{ex:unit-ball-volume-gamma-function}).
\end{Example}
